\documentclass[11pt]{article}
\usepackage[margin=1in]{geometry}
\usepackage{hyperref}
\usepackage{booktabs}
\usepackage{longtable}
\usepackage{xcolor}
\usepackage{graphicx}

\title{Replication Report: Tekedar et al.\ (2019)\\
\large ``Comparative genomics of \emph{Aeromonas veronii}: Identification of a pathotype impacting aquaculture globally''}
\author{Ollie (OpenClaw AI) --- BV-BRC Replication Wave}
\date{2026-07-01}

\begin{document}
\maketitle

\section*{Metadata}
\begin{tabular}{ll}
\toprule
Paper & Tekedar HC \emph{et al.}, PLoS ONE 14(9):e0221018 (2019) \\
DOI & \href{https://doi.org/10.1371/journal.pone.0221018}{10.1371/journal.pone.0221018} \\
PMC / PMID & PMC6715197 / 31465454 \\
License & CC BY 4.0 (open access) \\
Set & BVBRC-34 (TOPUP85 rank-15, $\sim$67 citations) \\
Compute & uicgpu (8$\times$A100), free open-source tools only \\
\textbf{Verdict} & \textbf{PARTIAL REPLICATION (strong)} \\
\bottomrule
\end{tabular}

\section{Paper summary}
The paper sequences and closes the genome of \emph{A.\ veronii} \textbf{ML09-123}, isolated
from a 2009 outbreak of motile \emph{Aeromonas} septicemia in farm-raised catfish
in the southeastern USA, and compares it against \textbf{all 40 other publicly
available \emph{A.\ veronii} genomes on NCBI as of 2/21/2018} (41 total; Table 1 of the paper),
spanning human, cattle, fish, water, and sediment sources across nine countries.

The paper's key conclusions:
\begin{enumerate}
  \item ML09-123 is highly similar to the Chinese isolate \textbf{TH0426}, suggesting a
        common origin and a shared aquaculture pathotype in the US and China
        (\emph{title claim}).
  \item All 41 genomes form a coherent species set; $\sim$4.3--5.0\,Mb, $\sim$58\% GC.
  \item Pan-genome $= 8710$, core-genome $= 2855$ (extrapolated 2791); core $\approx 30.9\%$
        of the pan; pan-genome remains \textbf{open}.
  \item Secretion systems T1SS, T2SS, T4P, and flagellum are conserved in \emph{all};
        T3SS, T5SS, T6SS, and TAD are variably distributed.
  \item $\sim$30\% of genomes show considerable variation in putative virulence genes.
\end{enumerate}

\section{Claims and testability}
\begin{longtable}{clll}
\toprule
\# & Claim & Testable? & Tested here? \\
\midrule
C1 & ML09-123 $\approx$ TH0426 (aquaculture pathotype) & yes & yes (3 methods) \\
C2 & 41/41 genuine \emph{A.\ veronii}; size + GC in range & yes & yes (41/41) \\
C3 & Pan 8710, core 2855, $\sim$30.9\%, open pan & yes & yes (CD-HIT) \\
C4 & T1SS/T2SS/T4P/flagellum conserved; T3SS/T5SS/T6SS/TAD variable & yes & yes (7/8) \\
C5 & $\sim$30\% strain variation in virulence genes & yes & yes (qualitative) \\
\bottomrule
\end{longtable}

\section{Method (free-tool substitutions for the paper's pipeline)}
\begin{itemize}
  \item Corpus: Europe PMC OA XML + PDF; Table~1 accession list; NCBI Datasets v2 to
        resolve each 2018 WGS project or strain to a current assembly. 41/41 resolved.
  \item Genome stats: length, GC\%, contig count, protein count for each of the 41 assemblies.
  \item ANI: \textbf{fastANI} all-vs-all (1{,}681 pairs) on the 41 nucleotide genomes.
  \item Phylogeny: \textbf{mash} sketch (s=100000) $\to$ all-vs-all distance $\to$
        average-linkage dendrogram.
  \item Pan/core: 166{,}630 predicted proteins clustered with \textbf{CD-HIT} at
        70\%/70\% (identity/coverage) --- a standard orthology proxy for EDGAR's
        BLAST-score-ratio criterion.
  \item Virulence factors: \textbf{abricate 1.4.0} $+$ bundled \textbf{VFDB}
        (4{,}592 sequences, the same database the paper used); strain $\times$ gene
        presence matrix; secretion-system / T4P / flagellum / TAD keyword buckets.
  \item Independent LLM judge: \texttt{argo:gpt-5.2} (free) --- opus-4.8 returned a
        transient 502; fell back per the wave brief's free-endpoint rule.
\end{itemize}

\section{Results vs paper}

\subsection*{C1 --- ML09-123 $\approx$ TH0426 (title claim) \textcolor{green!60!black}{REPRODUCED $\times$3}}
\begin{tabular}{lll}
\toprule
Method & Result & Interpretation \\
\midrule
fastANI ML09-123 $\times$ TH0426 & \textbf{99.927\%} & essentially clonal \\
fastANI 2nd-nearest to ML09-123 & AVNIH2 / CCM4359 $=$ 96.484\% & TH0426 is a clear outlier match \\
mash nearest neighbour of ML09-123 & \textbf{TH0426, $d = 0.00171$} & next-nearest Hm21 $= 0.03151$ ($\sim$18$\times$ farther) \\
VFDB profile & 136 vs 136 VFs, shared 136 & \textbf{Jaccard $= 1.000$} (identical) \\
\bottomrule
\end{tabular}

Three orthogonal methods (whole-genome ANI, k-mer distance, virulence-gene content)
independently confirm the paper's central pathotype claim.

\subsection*{C2 --- Species coherence + genome statistics \textcolor{green!60!black}{REPRODUCED (41/41)}}
All 41 genome sizes and GC\% match Table~1 (size delta $< 0.15$\,Mb; GC delta $< 0.6$\%).
Hm21 differs by 0.082\,Mb only because NCBI upgraded it to a complete assembly
(GCF\_000464515.\textbf{2}) since 2018. All 1{,}640 non-self fastANI pairs fall in
$[95.9\%, 100.0\%]$ (mean 96.9\%); \textbf{0 pairs} below the 95\% species threshold.

\subsection*{C3 --- Pan/core genome \textcolor{orange}{PARTIALLY REPRODUCED}}
\begin{tabular}{lrrl}
\toprule
Quantity & Paper (EDGAR SRV) & This work (CD-HIT 70/70) & Match \\
\midrule
Core (all/$\sim$all) & \textbf{2855} (extrap.\ 2791) & \textbf{2834} & \textcolor{green!60!black}{within 0.7\%} \\
Core (strict 100\% of 41) & --- & 1780 & (stricter cut) \\
Pan-genome & \textbf{8710} & 9664 & \textcolor{orange}{$+11\%$ (algorithm)} \\
Core \% of pan & \textbf{30.9\%} & \textbf{29.3\%} & \textcolor{green!60!black}{near-exact} \\
\bottomrule
\end{tabular}

Core size and core-fraction reproduce closely; the pan-genome total is $\sim$11\%
higher because CD-HIT's greedy 70\% clustering splits divergent paralogs/singletons
that EDGAR's BLAST-score-ratio criterion may merge. The gene-frequency spectrum
(3{,}319 genome-unique ``cloud'' genes) independently supports the paper's
\emph{open pan-genome} conclusion.

\subsection*{C4 --- Secretion-system conservation/variability \textcolor{green!60!black}{REPRODUCED (T5SS caveat)}}
\begin{tabular}{lrll}
\toprule
System & \# genomes (of 41) & Paper & Reproduced? \\
\midrule
T1SS       & 41 & conserved & \textcolor{green!60!black}{yes} \\
T2SS       & 41 & conserved & \textcolor{green!60!black}{yes} \\
T4P        & 41 & conserved & \textcolor{green!60!black}{yes} \\
Flagellum  & 41 & conserved & \textcolor{green!60!black}{yes} \\
T3SS       & 31 & variable  & \textcolor{green!60!black}{yes} \\
T6SS       & 28 & variable  & \textcolor{green!60!black}{yes} \\
TAD        & 12 & variable  & \textcolor{green!60!black}{yes} \\
T5SS       & 0 (unlabelled) & variable & \textcolor{orange}{un-testable in VFDB} \\
\bottomrule
\end{tabular}

7/8 systems reproduce the paper's exact conservation/variability call. T5SS is
un-testable because autotransporters are not annotated as ``type V secretion'' in
the VFDB product labels used.

\subsection*{C5 --- Substantial virulence-gene variation \textcolor{green!60!black}{QUALITATIVELY REPRODUCED}}
159 distinct VFDB genes across the 41 genomes; 58 core (in all 41), 101 variable (63.5\%).
Per-genome VF load: min 66, mean 117.5, max 140. Direction and biology agree with the
paper; the exact ``$\sim$30\% of genomes'' figure differs because abricate/VFDB and the
paper's CLC Genomic Workbench $+$ VFDB-setB screen use different DB scope and thresholds.

\section{Coverage / agreement}
Coverage: \textbf{5/5} claims meaningfully tested on real public data with real tool
runs. Agreement: \textbf{5/5} tested claims agree with the paper in direction, and
in magnitude for C1, C2, C3-core, and C4. \textbf{No claim contradicted.}
Independent LLM judge (\texttt{argo:gpt-5.2}, free): PARTIAL, coverage 5/5, agreement 5/5.

\section{Genuine critique}
This is a strong PARTIAL, not a full REPLICATED. I want to be explicit about what
did \emph{not} replicate cleanly, why, and what would be needed to close the gaps.

\begin{enumerate}
  \item \textbf{Pan-genome total (8710 vs 9664, $+11\%$).} This is a genuine
        \emph{algorithmic} difference, not a bug: CD-HIT at 70/70 is a greedy
        similarity-clusterer, whereas EDGAR 2.0 uses a bidirectional BLAST
        score-ratio (BSR) with an implicit merge threshold that groups distant
        paralogs the CD-HIT step will split. Reporting the two counts as ``the
        same experiment'' would be dishonest --- they are cousin methods, not the
        same method. Closing this gap would require re-running EDGAR itself
        (or Roary/Panaroo with matched thresholds), which was out of scope for
        the free-tool wave.
  \item \textbf{T5SS un-testable.} The abricate/VFDB screen used here does not
        label the autotransporter family as ``type V secretion,'' so T5SS is
        neither confirmed nor refuted --- an annotation-scope gap, not evidence
        against the paper. A proper T5SS assay would use a dedicated
        autotransporter HMM screen (SecReT5, TXSScan\_models); this is honestly
        flagged rather than papered over.
  \item \textbf{``$\sim$30\% of strains show substantial variation'' (C5).} I reproduce
        the \emph{direction} (wide per-genome VF-load spread; 63.5\% of VF genes are
        accessory) but the paper's exact ``$\sim$30\% of genomes'' framing is a
        threshold-and-DB choice (CLC Workbench + VFDB-setB, E$<$1e-50) that I
        cannot reproduce numerically without re-running the paper's specific
        pipeline. The qualitative agreement is real; the exact percentage
        agreement is not claimed.
  \item \textbf{Phylogeny is mash-distance based, not MUSCLE + NJ on the core
        alignment.} Mash is an excellent \emph{proxy} for core-genome distance
        but it is not the paper's phylogenetic method. The ML09-123 $\leftrightarrow$ TH0426
        clade is robustly recovered nonetheless, and fastANI + identical VFDB
        profile agree, so the C1 conclusion is well supported --- but the tree
        topology beyond that clade should not be over-interpreted.
  \item \textbf{Un-testable / not tested here:} host-of-isolation clustering
        (whether clinical vs aquaculture vs environmental \emph{A.\ veronii}
        cleanly separate by pathotype); temporal spread of the aquaculture
        pathotype since 2018 (1{,}927 genomes now vs 41 then --- a >45$\times$
        expansion that would strengthen or weaken the ``globally disseminated''
        framing); plasmid vs chromosomal location of the variable T3SS/T6SS/TAD
        loci; MDR / AMR carriage. These are genuinely open questions that the
        41-genome 2018 study did not resolve and this replication did not attempt.
  \item \textbf{Wave-scope honesty.} $\sim$30\,min of compute on a single A100 node
        with free tools cannot substitute for the paper's full comparative-genomics
        pipeline. The right claim is: the paper's \emph{conclusions} are
        independently supported on real public data; the paper's exact
        \emph{numbers} are reproduced where free tools permit and honestly
        differ where they do not.
\end{enumerate}

\section{Verdict}
\textbf{PARTIAL REPLICATION (strong).} The central pathotype claim (C1) is
reproduced by three independent methods; species coherence (C2) and the
secretion-system pattern (C4) are reproduced; the core-genome size (C3) matches
within 0.7\%; the pan-genome total differs by algorithm ($+11\%$); C5 agrees in
direction. No claim is contradicted.

\section{Resources (all free)}
Europe PMC REST; NCBI Datasets v2; fastANI; mash; CD-HIT; abricate 1.4.0 $+$
bundled VFDB; Argo proxy (\texttt{argo:gpt-5.2}); uicgpu 8$\times$A100 node
($\sim$30\,min wall).

\end{document}
